% SOURCE_AUTHORITY: sga5_fr_workpass.tex, lines 13996--14009 (LF-normalized unit).
% SOURCE_SHA256: 791F4EFFC5E02832D5D77ED03518C8156D6F07E4C8238B03545DB93D883FBB28
\subsection*{6.4. Observações sobre os termos locais}
\subsubsection*{a)} É fácil ver que, se $g\in G$, tem-se a relação:

\begin{equation}\tag{6.4.1}
\alpha_x(f,u,g^{-1}f',G,\varphi_g)=\alpha_x(f,u,f',G,\varphi)
\end{equation}
onde se pôs

\begin{equation}\tag{6.4.2}
\varphi_g(\xi)=(\operatorname{int}_g\circ\varphi)(\xi)=g^{-1}\varphi(\xi)g,
\quad \forall\,\xi\in G.
\end{equation}

% SOURCE_SCAN_NORMALIZATION: printed p.433 has X'_x^f here; the frozen authority certifies X'_x^{f'}, consistent with f' acting on X'.
A relação 6.4.1 permite sempre reduzir, substituindo, se necessário, $f'$ por $g^{-1}f'$ e $\varphi$ por $\varphi_g$, ao caso em que existe pelo menos um ponto $x'\in X_x^{\prime f'}$.
